\section{Public Reproducible Cohort}

We additionally evaluate a public, fully rerunnable cohort that is deliberately
separate from the confirmatory protocol. It contains twelve frozen local
transfer fixtures spanning Python, TypeScript, and Go: four predecessor-policy
transfers, four current-source-sufficient controls, and four stale or
distracting-memory cases. Each case declares the editable source roots, a
public behavioral verifier, and a dependency card that records both useful and
potentially misleading predecessor evidence.

The public cohort fixes one provider-reported model, one bounded agent tool
surface, a 300-second trial timeout, a per-trial cost ceiling, and three
predeclared repetitions per condition. We compare no memory, Mem0,
AgentMemory, and canonical \memorix{} retrieval. The primary contrast is
canonical \memorix{} versus no memory. Each provider receives the registered
predecessor record through its documented local interface; we record formation,
retrieval, model, tool, patch, and verifier receipts for every row.

The analysis collapses repetitions within a case before inference, so the
statistical unit remains the case rather than the retry. It reports paired
success-rate differences, case-cluster bootstrap intervals, and descriptive
resource differences for wall time, tokens, provider cost, and tool calls. The
cohort has only twelve independent case clusters, so it has low sensitivity to
modest effects and its interval width is a limitation rather than indirect
validation of any memory system. The
public analysis is fail-closed on missing, duplicate, mixed-model, or invalid
rows. Its results are explicitly descriptive: public fixtures cannot establish
pretraining novelty, private-oracle integrity, or general performance across
agents and repositories. We report them to make the complete implementation
and its observed boundaries reproducible while reserving causal or confirmatory
claims for the sealed-worker protocol.
